\chapter{Introduction}\label{ch:intro}

\section{Background and Motivation}

Chess is no longer only a board game. It has become a large-scale testbed for artificial
intelligence, from hand-crafted alpha-beta searchers to self-taught reinforcement-learning systems
such as AlphaZero~\cite{silver2018general}. Online platforms record hundreds of millions of games
every year, and open databases such as the Lichess puzzle collection~\cite{lichessdb} turn that
activity into millions of rated training positions. For the first time, the way ordinary club
players make decisions can be studied statistically, at scale.

The most popular training tool on these platforms is the tactical puzzle: a position with a single
clearly best continuation, which the player must find. Research on expertise suggests that practice
is effective in proportion to how precisely it targets a learner's current deficiencies, and that
repeating skills that are already mastered adds little~\cite{ericsson1993role}. Puzzle platforms,
however, mostly select puzzles by \emph{rating}. A player who handles forks well but repeatedly
misses back-rank mates still receives mostly forks, because forks are common. The practice is
comfortable, and largely redundant.

The original motivation for this project, as set out in the portfolio, was more ambitious. It was to
recognise a player's style of play (positional, attacking, defensive) from their move sequences,
and to generate new puzzles for that player based on the mistakes in their recent games. Over the
course of the project the emphasis shifted, in response to evidence gathered along the way, from
\emph{style} and \emph{generation} towards \emph{diagnosis} and \emph{adaptive selection}: finding
out what a player gets wrong, and serving the right puzzle at the right difficulty.
Section~\ref{sec:scope} documents this shift explicitly.

\section{Problem Statement}

Given a player's recent game history and a large pool of candidate puzzles, choose a sequence of
puzzles that maximises training value for that player. Four difficulties make this hard:

\begin{enumerate}
  \item \textbf{Weakness is not directly observable.} It must be inferred from noisy binary
        evidence (a puzzle is solved or not; a critical move in a game is found or missed), and every
        observation costs the learner time.
  \item \textbf{Tactic labels are noisy.} Deciding automatically \emph{which} motif a missed move
        belongs to is itself an unsolved classification problem.
  \item \textbf{Difficulty scales do not match.} A player's game rating on Chess.com and a puzzle's
        rating on Lichess are different scales, so ``puzzles near your rating'' is poorly defined.
  \item \textbf{Data per player is scarce.} A typical analysis covers about 50 games and roughly 80
        critical positions, spread across more than a dozen tactical categories.
\end{enumerate}

\section{Aim and Research Objectives}

The aim stated in the portfolio was \emph{``to design and evaluate an adaptive puzzle generator for
chess puzzles that helps players to enhance their skills by using reinforcement learning and player
style classification.''} The four objectives, and their status at the time of writing, are listed
in Table~\ref{tab:objectives}.

\begin{table}[H]
\centering
\small
\caption{Portfolio objectives and their status.}\label{tab:objectives}
\begin{tabularx}{\linewidth}{p{0.3\linewidth}Yp{0.13\linewidth}}
\toprule
\textbf{Objective (portfolio wording)} & \textbf{What was done} & \textbf{Status} \\
\midrule
O1. A classification model that predicts player rating group or style from gameplay data, with
statistically significant performance on unseen players. &
Rating-band classifier on 300 real players, cross-validated over players: 59.1\% accuracy over five
bands (chance 20\%), permutation $p=0.002$ (Section~\ref{sec:rating-classifier}). Style is covered
only by a rule-based archetype heuristic. & Met (rating group); style not met \\
O2. An adaptive puzzle generator that creates legal, unique and novel puzzles, optimised for
counter-intuitiveness and realism. &
Engine-verified \emph{mining} of puzzles from the player's own games through four Stockfish gates
(legal and unique by construction), plus a 5.9M-puzzle Lichess pool. No generative model, and no
counter-intuitiveness optimisation. & Partially met (reframed) \\
O3. A difficulty-matching module that selects puzzles according to player strength and prior
interaction history. &
An online Bayesian IRT learner that pitches each puzzle at a predicted 65\% solve probability, per
category, updated after every attempt (Section~\ref{sec:rl-setup}). & Met \\
O4. Evaluation with automatic puzzle-quality metrics and human assessment, showing improved fit
compared with a non-adaptive baseline. &
Automatic metrics, a pre-registered replication, a paired simulation study (seven experiments) and a
300-player real-data evaluation. The human study is designed but has not been run. & Partially met \\
\bottomrule
\end{tabularx}
\end{table}

\section{Research Questions}

The objectives lead to five research questions, each answered in Section~\ref{sec:answers}:

\begin{description}
  \item[RQ1] Can a player's rating group be predicted from \emph{how} they play, on players the model
        has never seen, without access to any rating information? (O1)
  \item[RQ2] Can a player's tactical weaknesses be diagnosed from their game history reliably
        enough to personalise training? (O1, O3)
  \item[RQ3] How reliably can the tactical motif of a position be labelled automatically, and does
        label quality matter downstream? (O2)
  \item[RQ4] Does difficulty-aware adaptive selection match puzzles to the player better than a
        fixed rating band? (O3)
  \item[RQ5] Does adaptive selection outperform non-adaptive baselines, and under which
        assumptions? (O4)
\end{description}

\section{Scope and Assumptions}\label{sec:scope}

\begin{honest}[Honest status: deviations from the portfolio plan]
\begin{itemize}
  \item \textbf{No generative model or RL fine-tuning of puzzles was built.} The portfolio's
        research gap, ``style- and rating-aware puzzle generation from scratch using generative RL
        conditioned on the learner'', remains open. The comparable published system (Feng et
        al.~\cite{feng2025generating}) pre-trains on 4.4 million puzzles and fine-tunes with RL at
        industrial-lab scale. Puzzles here are \emph{mined} from real games and verified by engine,
        so ``novel'' means novel \emph{to the player}, not new to the world.
  \item \textbf{Style classification was not built as a machine-learning model.} Objective O1 is met
        through the alternative the portfolio allowed, the \emph{rating group}.
  \item \textbf{No human assessment has yet been carried out.} Chapter~\ref{ch:discussion} and
        Appendix~\ref{app:repro} refer to a pre-registered protocol with a power analysis (34 paired
        participants are needed to detect a medium effect).
\end{itemize}
\end{honest}

\noindent\textbf{Assumptions.}
\begin{itemize}
  \item Stockfish~\cite{stockfish} at the node budgets used here (75,000 nodes per position for game
        analysis, depth 18 for puzzle verification) is treated as ground truth for move quality.
  \item Lichess puzzle themes are treated as a \emph{reference} for tactic labels. They are produced
        by an automatic tagger themselves, so agreement is not the same as accuracy.
  \item The probability of solving a puzzle follows a logistic (Rasch-type) form in the gap between
        ability and difficulty, as Elo and Glicko also assume~\cite{elo1978rating,glickman1999parameter}.
  \item Publicly available Chess.com games may be analysed for research when they are pseudonymised
        and reported only in aggregate (see the Declaration).
\end{itemize}

\section{Contributions of the Dissertation}

\begin{enumerate}
  \item \textbf{A working end-to-end system.} It covers Chess.com ingestion, deterministic Stockfish
        analysis graded by win probability, critical-position (opportunity) logging, weakness
        profiling, puzzle mining and adaptive serving, delivered as a Flask web application with 231
        offline tests.
  \item \textbf{A difficulty-aware Bayesian recommender.} An online item-response learner with a
        full-covariance Laplace update that reduces exactly to Glicko's update in one dimension
        (tested to $10^{-9}$). It combines Thompson Sampling over per-category ability offsets with
        per-puzzle difficulty targeting.
  \item \textbf{A validated line-level tactic tagger.} On a held-out sample, agreement with Lichess
        themes rises from 31.2\% ($\kappa=0.23$) to 58.4\% ($\kappa=0.52$).
  \item \textbf{PuzzleNet, a multi-task neural model of tactic and difficulty.} A network trained on
        5.5M Lichess puzzles that reads a position and the line that follows it and predicts the
        tactical category, the theme tags, and the difficulty \emph{with an uncertainty}. The rating
        likelihood treats each puzzle's Lichess rating deviation as known measurement noise, so the
        learned variance carries only what the content leaves uncertain, and that variance is then
        consumed directly by the Bayesian recommender as a puzzle's difficulty uncertainty. On the same
        held-out puzzles as the rule-based tagger, agreement rises to \AGREENN{} ($\kappa=\KAPPANN{}$), and
        mined puzzles gain a calibrated difficulty estimate in place of a closed-form guess. The network,
        its gradients and its optimiser are implemented from first principles in NumPy and checked
        against finite differences.
  \item \textbf{A rigorous comparative evaluation.} A pre-registered replication (reported
        including its failures), seven paired simulation experiments with Holm-corrected tests, and a
        prequential calibration on real solve logs.
  \item \textbf{A 300-player real-data evaluation.} Rating group is predictable from play; per-category
        weakness is almost entirely noise at typical data volumes; the rule-based and
        synthetic-trained profilers carry no predictive information.
  \item \textbf{Population category norms} fitted on the cohort. They separate ``this motif is hard for
        everyone'' from ``this player is weak at it''.
  \item \textbf{Documented engineering defects and fixes.} Several defects silently disabled or
        corrupted parts of the system; they are reported in Section~\ref{sec:challenges} because they
        bear on the validity of earlier results.
\end{enumerate}

\section{Dissertation Outline}

Chapter~\ref{ch:literature} reviews computer chess, puzzle platforms, player modelling, generative
models and RL, and difficulty modelling. Chapter~\ref{ch:problem} formalises the problem and presents
the system architecture. Chapter~\ref{ch:method} describes the methodology, and
Chapter~\ref{ch:implementation} the implementation. Chapter~\ref{ch:results} reports the experiments.
Chapter~\ref{ch:discussion} discusses the findings, including an honest assessment of how
sophisticated the AI in this system is. Chapter~\ref{ch:conclusion} answers the research questions
and recommends future work.
